% Mathématiques 6e — cours V3
% Source LaTeX rendue côté Astro/KaTeX.

\section{Objectifs}

Dans ce chapitre, on apprend à :

\begin{itemize}
\item raisonner sur une quantité dont on ne connaît pas encore la valeur ;
\item utiliser des schémas, des symboles ou des lettres comme outils de représentation ;
\item comprendre le sens d’une égalité ;
\item chercher une valeur inconnue sans appliquer mécaniquement une technique d’équation ;
\item lire, exécuter et produire des programmes de calcul simples ;
\item reconnaître une régularité dans une suite de nombres ou de motifs ;
\item prévoir une étape éloignée d’un motif évolutif ;
\item expliquer et justifier une règle observée.
\end{itemize}

\section{1. La pensée algébrique commence avant le calcul littéral}

En 6e, l’algèbre ne consiste pas à apprendre une longue liste de règles sur les lettres.

L’objectif est d’apprendre à raisonner sur des \textbf{relations entre des quantités}, parfois sans connaître immédiatement leur valeur.

Par exemple, dans la situation :

> Deux objets identiques et un objet de valeur $3$ ont ensemble une valeur de $17$.

on peut représenter la valeur inconnue d’un objet par un symbole.

Si cette valeur est notée $x$, la situation peut être traduite par :

\[
x+x+3=17.
\]

Mais la lettre n’est pas le point de départ obligatoire : un schéma en barres ou une représentation concrète peut être plus efficace.

\section{2. Le signe égal exprime une égalité de valeurs}

Le signe

\[
=
\]

signifie que les deux expressions placées de part et d’autre ont \textbf{la même valeur}.

Par exemple :

\[
7+5=3\times4.
\]

Cette égalité est vraie car :

\[
7+5=12
\]

et :

\[
3\times4=12.
\]

Le signe égal ne signifie donc pas seulement « maintenant, je donne le résultat ».

Il peut relier deux expressions différentes qui représentent le même nombre.

\section{3. Égalités vraies et fausses}

Considérons :

\[
8+6=7\times2.
\]

Les deux membres valent :

\[
14.
\]

L’égalité est donc vraie.

En revanche :

\[
8+6=5\times3
\]

est fausse car :

\[
14\neq15.
\]

Pour vérifier une égalité, on calcule les deux membres \textbf{séparément} puis on compare leurs valeurs.

\section{4. Représenter une quantité inconnue}

Une quantité inconnue peut être représentée :

\begin{itemize}
\item par une case vide ;
\item par un symbole ;
\item par une barre ;
\item éventuellement par une lettre.
\end{itemize}

Par exemple :

\[
\Box+9=21
\]

ou :

\[
x+9=21
\]

traduisent la même recherche : trouver un nombre qui, augmenté de $9$, donne $21$.

Le nombre cherché est :

\[
21-9=12.
\]

Donc :

\[
x=12.
\]

\section{5. Modèle en barres}

Les schémas en barres permettent de visualiser les relations entre les quantités.

\coursefigure{/images/mathematiques/6eme/algebre-modele-barres.svg}{Schéma en barres représentant une quantité inconnue complétée par une quantité connue pour former une quantité totale.}{Si une quantité inconnue $x$ augmentée de $6$ donne $20$, le schéma représente la relation $x+6=20$.}

On lit alors directement :

\[
x=20-6.
\]

Donc :

\[
\boxed{x=14}.
\]

Le schéma aide à identifier l’opération à effectuer sans réciter une règle abstraite.

\section{6. Exemple développé 1 — retrouver une quantité inconnue}

Trois boîtes identiques et $4$ objets isolés contiennent ensemble $25$ objets.

Chaque boîte contient le même nombre $x$ d’objets.

On peut écrire :

\[
3x+4=25.
\]

On commence par retirer les $4$ objets isolés :

\[
25-4=21.
\]

Les trois boîtes identiques contiennent donc ensemble $21$ objets.

Chaque boîte contient :

\[
21\div3=7.
\]

Ainsi :

\[
\boxed{x=7}.
\]

Vérification :

\[
3\times7+4=25.
\]

Cette démarche repose sur le sens des opérations.

\section{7. Opérations inverses}

Les opérations inverses sont utiles pour retrouver une quantité inconnue.

L’addition et la soustraction sont liées.

Par exemple :

\[
x+8=31
\]

conduit à :

\[
x=31-8=23.
\]

La multiplication et la division sont liées.

Par exemple :

\[
5x=35
\]

conduit à :

\[
x=35\div5=7.
\]

En 6e, on privilégie ce raisonnement arithmétique avant les techniques formelles d’équations du cycle suivant.

\section{8. Une lettre peut désigner un nombre variable}

Une lettre ne représente pas toujours un nombre que l’on cherche.

Elle peut aussi représenter un nombre qui \textbf{peut prendre plusieurs valeurs}.

Supposons qu’un objet coûte $12\ \text{euros}$ et qu’un forfait fixe de livraison coûte $5\ \text{euros}$.

Si $n$ désigne le nombre d’objets achetés, le prix total est :

\[
12n+5.
\]

Pour $n=3$ :

\[
12\times3+5=41.
\]

Pour $n=8$ :

\[
12\times8+5=101.
\]

La même expression décrit donc plusieurs situations.

\section{9. Substitution}

Substituer une valeur à une lettre signifie remplacer la lettre par cette valeur.

Considérons :

\[
4n+3.
\]

Pour :

\[
n=5,
\]

on calcule :

\[
4\times5+3=23.
\]

Pour :

\[
n=10,
\]

on calcule :

\[
4\times10+3=43.
\]

La substitution sert à tester ou exploiter une règle.

\section{10. Programmes de calcul}

Un programme de calcul est une suite ordonnée d’instructions.

Exemple :

\begin{enumerate}
\item choisir un nombre ;
\item le multiplier par $3$ ;
\item ajouter $5$.
\end{enumerate}

Si le nombre choisi est $4$ :

\[
4\times3=12,
\]

puis :

\[
12+5=17.
\]

Le résultat final est :

\[
17.
\]

\section{11. Traduire un programme de calcul}

Si le nombre de départ est noté $n$, le programme précédent donne :

\[
3n+5.
\]

Cette écriture permet de décrire le programme \textbf{sans choisir une valeur particulière}.

Attention : en 6e, l’objectif n’est pas encore de développer tout le calcul littéral. La lettre sert surtout à représenter une quantité de manière générale.

\section{12. Exécuter un programme à l’envers}

On peut parfois retrouver le nombre de départ en remontant les opérations.

Programme :

\begin{enumerate}
\item choisir un nombre ;
\item ajouter $7$ ;
\item multiplier le résultat par $2$.
\end{enumerate}

Le résultat final est $30$.

On remonte le programme :

\[
30\div2=15,
\]

puis :

\[
15-7=8.
\]

Le nombre choisi était donc :

\[
\boxed{8}.
\]

\section{13. Comparer deux programmes}

Programme A :

\begin{enumerate}
\item choisir un nombre ;
\item le doubler ;
\item ajouter $6$.
\end{enumerate}

Programme B :

\begin{enumerate}
\item choisir un nombre ;
\item ajouter $3$ ;
\item doubler.
\end{enumerate}

Pour un nombre $n$, le programme A donne :

\[
2n+6.
\]

Le programme B donne :

\[
2(n+3).
\]

Or :

\[
2(n+3)=2n+6.
\]

Les deux programmes donnent donc toujours le même résultat.

Dans ce cas, quelques essais peuvent suggérer l’égalité, mais l’écriture générale explique \textbf{pourquoi} elle est toujours vraie.

\section{14. Essais et conjectures}

Tester une règle sur quelques exemples permet de formuler une \textbf{conjecture}.

Une conjecture est une propriété que l’on pense vraie mais qui demande encore une justification.

Par exemple, on observe :

\[
1+3=4,
\]

\[
3+5=8,
\]

\[
5+7=12.
\]

On peut conjecturer :

> La somme de deux nombres impairs consécutifs est paire.

Les exemples rendent la conjecture plausible.

Mais un raisonnement général est nécessaire pour expliquer qu’elle ne dépend pas seulement des exemples choisis.

\section{15. Motifs évolutifs}

Un motif évolutif change d’une étape à la suivante selon une règle.

\coursefigure{/images/mathematiques/6eme/algebre-motif-evolutif.svg}{Quatre étapes d’un motif évolutif construit avec un nombre croissant de carrés.}{Dans cet exemple visuel, une nouvelle case est ajoutée à chaque étape. L’objectif est d’identifier la structure, puis de prévoir une étape éloignée.}

La première question est :

> Qu’est-ce qui reste identique et qu’est-ce qui change ?

La seconde est :

> Comment décrire la règle sans avoir besoin de dessiner toutes les étapes précédentes ?

\section{16. Décrire une régularité}

Supposons qu’un motif comporte :

\[
4
\]

éléments à l’étape $1$, puis gagne :

\[
3
\]

éléments à chaque nouvelle étape.

Les premières valeurs sont :

\[
4,\quad7,\quad10,\quad13,\quad16,\ldots
\]

On peut dire :

> À chaque étape, on ajoute $3$.

Pour passer de l’étape $1$ à l’étape $n$, on effectue :

\[
n-1
\]

augmentations.

Le nombre d’éléments vaut donc :

\[
4+3(n-1).
\]

Cette écriture est une manière compacte de décrire la structure du motif.

\section{17. Exemple développé 2 — prévoir une étape éloignée}

Un motif contient $5$ éléments à l’étape $1$ et gagne $4$ éléments à chaque étape.

Combien contient-il d’éléments à l’étape $8$ ?

Entre les étapes $1$ et $8$, il y a :

\[
8-1=7
\]

augmentations.

On ajoute donc :

\[
7\times4=28
\]

éléments.

Le total vaut :

\[
5+28=33.
\]

Ainsi, l’étape $8$ contient :

\[
\boxed{33}
\]

éléments.

\section{18. Chercher une structure plutôt qu’une simple liste}

Une suite peut être décrite de plusieurs manières.

Considérons :

\[
2,\quad5,\quad8,\quad11,\quad14,\ldots
\]

On peut remarquer :

> On ajoute $3$ à chaque fois.

Mais on peut aussi relier directement le rang $n$ au nombre obtenu.

À l’étape $1$, le résultat est $2$.

Une expression possible est :

\[
2+3(n-1).
\]

Cette deuxième description permet d’atteindre directement une étape éloignée.

\section{19. Motifs non numériques}

La pensée algébrique ne concerne pas uniquement les suites de nombres.

Un motif géométrique peut aussi évoluer.

Par exemple :

\begin{itemize}
\item étape $1$ : un carré ;
\item étape $2$ : trois carrés ;
\item étape $3$ : cinq carrés ;
\item étape $4$ : sept carrés.
\end{itemize}

Les nombres d’éléments sont :

\[
1,\quad3,\quad5,\quad7,\ldots
\]

La régularité est :

> Ajouter $2$ à chaque étape.

On peut chercher combien d’éléments possède une étape éloignée sans construire toutes les étapes précédentes.

\section{20. Lire un tableau de valeurs}

Un tableau peut aider à observer une relation.

| Étape | Nombre d’éléments |
|---:|---:|
| $1$ | $4$ |
| $2$ | $7$ |
| $3$ | $10$ |
| $4$ | $13$ |

On observe que lorsque l’étape augmente de :

\[
1,
\]

le nombre d’éléments augmente de :

\[
3.
\]

Le tableau ne donne pas encore une preuve, mais il rend la régularité plus visible.

\section{21. Méthode — résoudre un problème pré-algébrique}

Lorsqu’un nombre inconnu intervient :

\subsection{Étape 1 — identifier les quantités}

Qu’est-ce qui est connu ?

Qu’est-ce qui est inconnu ?

\subsection{Étape 2 — représenter la relation}

On peut utiliser :

\begin{itemize}
\item un schéma en barres ;
\item une case vide ;
\item une lettre ;
\item une égalité.
\end{itemize}

\subsection{Étape 3 — raisonner avec les opérations}

On utilise les opérations inverses ou la structure de la situation.

\subsection{Étape 4 — vérifier}

On remplace la quantité inconnue par la valeur trouvée dans la situation initiale.

\section{22. Méthode — étudier un motif évolutif}

\subsection{Étape 1 — observer plusieurs étapes}

Repérer précisément ce qui est ajouté ou modifié.

\subsection{Étape 2 — décrire la transformation}

Par exemple :

> Chaque nouvelle étape ajoute $3$ éléments.

\subsection{Étape 3 — compter le nombre de transformations}

Pour passer de l’étape $1$ à l’étape $n$ :

\[
n-1
\]

transformations sont effectuées.

\subsection{Étape 4 — tester la règle}

La règle doit retrouver les premières étapes connues.

\section{23. Erreurs fréquentes}

\subsection{Utiliser le signe égal comme une flèche}

Il est incorrect d’écrire :

\[
5+3=8\times2=16
\]

car :

\[
5+3=8
\]

mais :

\[
8\neq16.
\]

Il faut écrire les étapes séparément.

\subsection{Croire qu’une lettre désigne toujours la même valeur}

Dans une expression générale, $n$ peut prendre différentes valeurs selon la situation.

\subsection{Introduire trop vite des techniques formelles}

Pour :

\[
x+7=18,
\]

il suffit en 6e de comprendre que :

\[
x=18-7=11.
\]

\subsection{Confondre l’étape et le nombre d’ajouts}

De l’étape $1$ à l’étape $8$, il y a :

\[
7
\]

changements, pas $8$.

\subsection{Croire que quelques exemples prouvent une règle}

Des exemples peuvent permettre de conjecturer.

Ils ne suffisent pas toujours à justifier une affirmation générale.

\section{24. À retenir}

La pensée algébrique consiste d’abord à \textbf{raisonner sur des relations}.

Une quantité inconnue peut être représentée par un symbole ou une lettre.

Une égalité affirme que deux expressions ont la même valeur :

\[
\boxed{A=B}.
\]

Les modèles en barres permettent de rendre visibles les relations entre quantités.

Un motif évolutif doit être analysé en recherchant :

\begin{itemize}
\item ce qui change ;
\item de combien cela change ;
\item combien de fois la transformation est répétée.
\end{itemize}

Une règle efficace doit être testée sur les premières étapes et permettre de prévoir des étapes plus éloignées.

L’objectif n’est pas de faire du calcul littéral avancé, mais d’apprendre à \textbf{généraliser, représenter et justifier}.
